% Options for packages loaded elsewhere
\PassOptionsToPackage{unicode}{hyperref}
\PassOptionsToPackage{hyphens}{url}
\documentclass[aspectratio=169,ignorenonframetext]{beamer}
\usepackage{ma2003b-beamer}
\hypersetup{
  pdftitle={Cluster Analysis},
  pdfauthor={Juliho Castillo Colmenares},
  hidelinks,
  pdfcreator={LaTeX via pandoc}}

\title{\texorpdfstring{Cluster Analysis}{Cluster Analysis}}
\author{\texorpdfstring{Juliho Castillo
Colmenares}{Juliho Castillo Colmenares}}
\subtitle{Theory, Methods, and Practical Application}
\institute{Tec de Monterrey}
\date{}

\begin{document}
\frame{\titlepage}

\begin{frame}[shrink]{Today's Agenda}
\textbf{Part 1: Theoretical Foundations}

\begin{enumerate}
\item
  Introduction to Cluster Analysis
\item
  Distance and Similarity Measures
\item
  Hierarchical Clustering Methods
\item
  K-Means and Non-Hierarchical Methods
\item
  Determining Optimal Number of Clusters
\item
  Validation Techniques
\item
  Practical Considerations
\item
  Applications and Best Practices
\end{enumerate}

\textbf{Part 2: Practical Application}

\begin{enumerate}
\item
  Customer Segmentation Case Study
\end{enumerate}
\end{frame}

\begin{frame}{Companion Case Study}
\textbf{E-Commerce Customer Segmentation}

Throughout this presentation, we'll reference a real-world example
\end{frame}

\begin{frame}{Case Study: Business Context}
\textbf{Scenario:} E-commerce company analyzing customer behavior

\textbf{Challenge:} Discover natural customer segments without
predefined categories

\textbf{Goal:} Develop targeted marketing strategies for each segment
\end{frame}

\begin{frame}{Case Study: Dataset}
\textbf{2,000 Customers} with 7 behavioral variables:

\begin{itemize}
\item
  monthly\_purchases
\item
  avg\_basket\_size
\item
  total\_spend
\item
  session\_duration
\item
  email\_clicks
\item
  product\_views
\item
  return\_rate
\end{itemize}

\textbf{No predefined labels} - pure unsupervised learning
\end{frame}

\begin{frame}{What We'll Discover}
\textbf{Spoiler Preview:}

\begin{itemize}
\item
  4 distinct customer segments
\item
  Each with unique behavioral patterns
\item
  Actionable marketing strategies
\end{itemize}

\textbf{You'll see how theory connects to practice throughout}
\end{frame}

\begin{frame}{What is Cluster Analysis?}
\textbf{Definition:} An exploratory technique to discover natural
groupings in data \textbf{without predefined categories}
\end{frame}

\begin{frame}{What is Cluster Analysis?}
\textbf{Key Characteristics:}

\begin{itemize}
\item
  Unsupervised learning method
\item
  No training labels required
\item
  Discovers hidden structure in data
\item
  Groups similar observations together
\end{itemize}

\textbf{Goal:} Maximize within-cluster similarity and between-cluster
dissimilarity
\end{frame}

\begin{frame}{Cluster Analysis vs. Discriminant Analysis}
{\def\LTcaptype{none} % do not increment counter
\begin{longtable}[]{@{}
  >{\raggedright\arraybackslash}p{(\linewidth - 2\tabcolsep) * \real{0.5000}}
  >{\raggedright\arraybackslash}p{(\linewidth - 2\tabcolsep) * \real{0.5000}}@{}}
\toprule\noalign{}
\endhead
\textbf{Cluster Analysis} & \textbf{Discriminant Analysis} \\
Unsupervised learning & Supervised learning \\
Discovers unknown groups & Classifies into known groups \\
No training labels & Requires training labels \\
Exploratory & Predictive \\
Groups observations & Creates decision boundaries \\
\bottomrule\noalign{}
\end{longtable}
}
\end{frame}

\begin{frame}{Applications: Marketing \& Business}
\textbf{Marketing}

\begin{itemize}
\item
  Customer segmentation for targeted campaigns
\item
  Market basket analysis
\end{itemize}

\textbf{Business}

\begin{itemize}
\item
  Fraud detection
\item
  Anomaly identification
\end{itemize}
\end{frame}

\begin{frame}{Applications: Science \& Healthcare}
\textbf{Biology \& Medicine}

\begin{itemize}
\item
  Disease subtype identification
\item
  Gene expression analysis
\end{itemize}

\textbf{Social Sciences}

\begin{itemize}
\item
  Community detection in networks
\item
  Document clustering
\end{itemize}
\end{frame}

\begin{frame}{Distance and Similarity Measures}
\textbf{Why Distance Matters}

Clustering depends on measuring how ``close'' observations are to each
other
\end{frame}

\begin{frame}{Common Distance Metrics}
\begin{enumerate}
\item
  \textbf{Euclidean Distance} (L2 norm) - Most common
\item
  \textbf{Manhattan Distance} (L1 norm) - Robust to outliers
\item
  \textbf{Cosine Similarity} - For high-dimensional data
\item
  \textbf{Correlation Distance} - Pattern similarity
\end{enumerate}
\end{frame}

\begin{frame}{Euclidean Distance}
\textbf{Formula:}

\[d(x,y) = \sqrt{\sum_{i = 1}^{p}\left( x_{i} - y_{i} \right)^{2}}\]
\end{frame}

\begin{frame}{Euclidean Distance}
\textbf{Properties:}

\begin{itemize}
\item
  Straight-line distance in n-dimensional space
\item
  Sensitive to scale differences
\item
  Assumes equal importance of all dimensions
\end{itemize}

\textbf{Warning:} Always standardize variables with different scales!
\end{frame}

\begin{frame}{Manhattan Distance}
\textbf{Formula:}

\[d(x,y) = \sum_{i = 1}^{p}|x_{i} - y_{i}|\]
\end{frame}

\begin{frame}{Manhattan Distance}
\textbf{When to Use:}

\begin{itemize}
\item
  Data contains outliers or extreme values
\item
  Variables represent counts
\item
  High-dimensional spaces
\end{itemize}

\textbf{Advantage:} More robust than Euclidean distance
\end{frame}

\begin{frame}{Why Standardization is Critical}
\textbf{Problem:} Variables on different scales dominate distance
calculations

\textbf{Example:}

\begin{itemize}
\item
  Age: 20-80 years
\item
  Income: 20,000-200,000 dollars
\end{itemize}

Without standardization, income dominates!
\end{frame}

\begin{frame}{Z-score Standardization}
\textbf{Solution: Z-score Standardization}

\[z_{i} = \frac{x_{i} - \mu}{\sigma}\]

Transform to mean = 0, standard deviation = 1
\end{frame}

\begin{frame}{Case Study: Why Standardization Matters}
\textbf{Customer Data Scale Problem:}

{\def\LTcaptype{none} % do not increment counter
\begin{longtable}[]{@{}
  >{\raggedright\arraybackslash}p{(\linewidth - 4\tabcolsep) * \real{0.5001}}
  >{\centering\arraybackslash}p{(\linewidth - 4\tabcolsep) * \real{0.2501}}
  >{\centering\arraybackslash}p{(\linewidth - 4\tabcolsep) * \real{0.2501}}@{}}
\toprule\noalign{}
\endhead
\textbf{Variable} & \textbf{Min} & \textbf{Max} \\
monthly\_purchases & 0.2 & 17.0 \\
total\_spend (dollars) & 58 & 7,892 \\
return\_rate (proportion) & 0.0 & 0.5 \\
\bottomrule\noalign{}
\end{longtable}
}

Without standardization, total\_spend would dominate clustering!
\end{frame}

\begin{frame}{Case Study: After Standardization}
\textbf{All variables transformed to mean ≈ 0, std ≈ 1}

Result: Each behavioral dimension contributes equally

\begin{itemize}
\item
  Monthly purchases: now comparable to spending
\item
  Allows discovery of patterns beyond just ``who spends most''
\item
  Reveals behavioral segments not visible in raw data
\end{itemize}
\end{frame}

\begin{frame}{Hierarchical Clustering}
Builds a tree-like structure (dendrogram) showing nested clusters
\end{frame}

\begin{frame}{Two Approaches}
\textbf{Agglomerative (Bottom-Up):} Most common

\begin{itemize}
\item
  Start: Each observation is its own cluster
\item
  Process: Merge closest clusters iteratively
\item
  End: All observations in one cluster
\end{itemize}
\end{frame}

\begin{frame}{Two Approaches}
\textbf{Divisive (Top-Down):} Less common

\begin{itemize}
\item
  Start: All observations in one cluster
\item
  Process: Split most heterogeneous cluster
\item
  End: Each observation is its own cluster
\end{itemize}
\end{frame}

\begin{frame}{Linkage Methods}
How to Measure Distance Between Clusters?
\end{frame}

\begin{frame}{Single Linkage (Nearest Neighbor)}
\[d\left( C_{1},C_{2} \right) = \min\limits_{x \in C_{1},y \in C_{2}}d(x,y)\]

Distance between closest points in the two clusters
\end{frame}

\begin{frame}{Complete Linkage (Farthest Neighbor)}
\[d\left( C_{1},C_{2} \right) = \max\limits_{x \in C_{1},y \in C_{2}}d(x,y)\]

Distance between farthest points in the two clusters
\end{frame}

\begin{frame}{Average Linkage}
\[d\left( C_{1},C_{2} \right) = \frac{1}{n_{1}n_{2}}\sum_{x \in C_{1}}\sum_{y \in C_{2}}d(x,y)\]

Average distance between all pairs of points
\end{frame}

\begin{frame}{Ward's Method}
\textbf{Ward's Method}

Minimizes within-cluster sum of squares

Tends to produce compact, equal-sized clusters
\end{frame}

\begin{frame}{Linkage Methods Comparison}
{\def\LTcaptype{none} % do not increment counter
\begin{longtable}[]{@{}
  >{\raggedright\arraybackslash}p{(\linewidth - 4\tabcolsep) * \real{0.3244}}
  >{\centering\arraybackslash}p{(\linewidth - 4\tabcolsep) * \real{0.2703}}
  >{\raggedright\arraybackslash}p{(\linewidth - 4\tabcolsep) * \real{0.4055}}@{}}
\toprule\noalign{}
\endhead
\textbf{Method} & \textbf{Outlier Sensitivity} & \textbf{Cluster
Shape} \\
Single Linkage & High & Elongated (chaining) \\
Complete Linkage & Low & Compact, spherical \\
Average Linkage & Medium & Balanced \\
Ward's Method & Medium & Compact, equal-sized \\
\bottomrule\noalign{}
\end{longtable}
}
\end{frame}

\begin{frame}{Linkage Methods Comparison}
\textbf{Recommendation:} Ward's method often works best in practice
\end{frame}

\begin{frame}{Dendrograms}
Visualizing Hierarchical Structure
\end{frame}

\begin{frame}{Reading a Dendrogram}
\begin{itemize}
\item
  Horizontal axis: Observations or clusters
\item
  Vertical axis: Distance at which clusters merge
\item
  Height of branches: Dissimilarity between merged clusters
\end{itemize}
\end{frame}

\begin{frame}{Determining Number of Clusters}
\begin{itemize}
\item
  Look for large vertical gaps (jumps in fusion distance)
\item
  Cut dendrogram where there's substantial increase
\item
  Draw horizontal line: number of vertical lines crossed = k clusters
\end{itemize}
\end{frame}

\begin{frame}{Case Study: Dendrogram Analysis}
\textbf{Ward's Method Results:}

\begin{itemize}
\item
  Clear hierarchical structure
\item
  Large vertical gap suggests 4 clusters
\item
  Cut at distance ≈ 50
\end{itemize}

\textbf{Cluster Sizes (Hierarchical):}

\begin{itemize}
\item
  Cluster 0: 440 customers (22\%)
\item
  Cluster 1: 300 customers (15\%)
\item
  Cluster 2: 560 customers (28\%)
\item
  Cluster 3: 700 customers (35\%)
\end{itemize}
\end{frame}

\begin{frame}{Case Study: Hierarchical Quality}
\textbf{Silhouette Score: 0.458}

\begin{itemize}
\item
  Above 0.4 is acceptable for customer segmentation
\item
  Behavioral boundaries are naturally fuzzy
\item
  Reasonably well-separated clusters
\item
  Balanced cluster sizes
\end{itemize}

Dendrogram suggests 4 is the natural number of segments
\end{frame}

\begin{frame}{The Chaining Effect}
\textbf{Problem with Single Linkage:}

Clusters form long, elongated chains rather than compact groups
\end{frame}

\begin{frame}{The Chaining Effect}
\textbf{Why it Happens:}

\begin{itemize}
\item
  Observations connect via intermediate points
\item
  A-B-C-D form chain where each is close to neighbor
\item
  But A and D are far apart
\end{itemize}
\end{frame}

\begin{frame}{The Chaining Effect}
\textbf{Solution:}

\begin{itemize}
\item
  Use complete or average linkage instead
\item
  Or Ward's method for compact clusters
\end{itemize}
\end{frame}

\begin{frame}{K-Means Clustering}
Most popular non-hierarchical method
\end{frame}

\begin{frame}{K-Means Algorithm}
\begin{enumerate}
\item
  \textbf{Initialize:} Select k random observations as centroids
\item
  \textbf{Assignment:} Assign each point to nearest centroid
\item
  \textbf{Update:} Recalculate centroids as cluster means
\item
  \textbf{Repeat:} Steps 2-3 until convergence
\end{enumerate}

\textbf{Convergence:} When assignments no longer change between
iterations
\end{frame}

\begin{frame}{K-Means Objective Function}
\textbf{Goal:} Minimize within-cluster sum of squares (WCSS)

\[\min\sum_{i = 1}^{k}\sum_{x \in C_{i}}\| x - \mu_{i}\|^{2}\]

where \(\mu_{i}\) is the centroid of cluster \(C_{i}\)
\end{frame}

\begin{frame}{K-Means Properties}
\begin{itemize}
\item
  Always converges (finite partitions, monotonically decreasing WCSS)
\item
  Typically converges in 10-30 iterations
\item
  Fast: O(n times k times p times iterations)
\end{itemize}
\end{frame}

\begin{frame}{K-Means: Advantages}
\begin{itemize}
\item
  Fast and scalable to large datasets
\item
  Simple to understand and implement
\item
  Efficient for exploratory analysis
\end{itemize}
\end{frame}

\begin{frame}{K-Means: Limitations}
\begin{itemize}
\item
  Requires specifying k in advance
\item
  Sensitive to initialization (different starts → different results)
\item
  Assumes spherical clusters
\item
  Sensitive to outliers
\item
  Tends to create equal-sized clusters
\end{itemize}
\end{frame}

\begin{frame}{K-Means++ Initialization}
\textbf{Problem:} Random initialization can lead to poor results
\end{frame}

\begin{frame}{K-Means++ Algorithm}
\begin{enumerate}
\item
  Choose first centroid randomly
\item
  For each subsequent centroid:

  \begin{itemize}
  \item
    Choose point with probability proportional to squared distance from
    nearest existing centroid
  \end{itemize}
\item
  Repeat until k centroids selected
\end{enumerate}

\textbf{Benefit:} Spreads out initial centroids, significantly improves
results
\end{frame}

\begin{frame}{K-Medoids (PAM)}
\textbf{Key Difference from K-Means:}

\begin{itemize}
\item
  K-means: Centers are computed means (may not be actual points)
\item
  K-medoids: Centers are actual data points (medoids)
\end{itemize}
\end{frame}

\begin{frame}{K-Medoids (PAM)}
\textbf{Advantages:}

\begin{itemize}
\item
  More robust to outliers
\item
  Works with any distance metric
\item
  Interpretable centers (actual observations)
\end{itemize}

\textbf{Disadvantage:} Slower than k-means (higher computational cost)
\end{frame}

\begin{frame}{How Many Clusters?}
\textbf{The Fundamental Challenge:}

No ``ground truth'' for correct number of clusters
\end{frame}

\begin{frame}{Multiple Approaches}
\begin{enumerate}
\item
  \textbf{Elbow Method} - Look for bend in WCSS plot
\item
  \textbf{Silhouette Analysis} - Measure cluster quality
\item
  \textbf{Gap Statistic} - Compare to null reference
\item
  \textbf{Davies-Bouldin Index} - Ratio of compactness to separation
\item
  \textbf{Domain Knowledge} - Business requirements
\end{enumerate}
\end{frame}

\begin{frame}{Elbow Method: Procedure}
\begin{enumerate}
\item
  Run clustering for k = 1, 2, 3, \ldots, K\_max
\item
  Calculate WCSS for each k
\item
  Plot WCSS vs. k
\item
  Look for ``elbow'' - diminishing returns point
\end{enumerate}
\end{frame}

\begin{frame}{Elbow Method: Interpretation}
\begin{itemize}
\item
  WCSS always decreases as k increases
\item
  Elbow indicates where additional clusters don't help much
\item
  Choose k at the elbow point
\end{itemize}

\textbf{Limitation:} Elbow not always clear - may need other methods
\end{frame}

\begin{frame}{Case Study: Elbow Method Results}
\textbf{Tested k = 2 to 10:}

\begin{itemize}
\item
  Elbow curve shows diminishing returns after k=4
\item
  Silhouette scores peak at k=4
\item
  Both metrics converge on same answer
\end{itemize}

\textbf{K-Means with k=4:}

\begin{itemize}
\item
  Silhouette Score: 0.458 (same as hierarchical!)
\item
  Convergence of two different methods
\item
  Strong evidence for 4 segments
\end{itemize}
\end{frame}

\begin{frame}{Case Study: K-Means Cluster Sizes}
\textbf{Final K-Means Segmentation (k=4):}

{\def\LTcaptype{none} % do not increment counter
\begin{longtable}[]{@{}
  >{\centering\arraybackslash}p{(\linewidth - 4\tabcolsep) * \real{0.3333}}
  >{\centering\arraybackslash}p{(\linewidth - 4\tabcolsep) * \real{0.3333}}
  >{\centering\arraybackslash}p{(\linewidth - 4\tabcolsep) * \real{0.3333}}@{}}
\toprule\noalign{}
\endhead
\textbf{Cluster} & \textbf{Size} & \textbf{Percentage} \\
0 & 307 & 15.3\% \\
1 & 707 & 35.4\% \\
2 & 432 & 21.6\% \\
3 & 554 & 27.7\% \\
\bottomrule\noalign{}
\end{longtable}
}

Hierarchical and K-means produce similar cluster sizes
\end{frame}

\begin{frame}{Silhouette Analysis}
Measures how well each point fits within its cluster
\end{frame}

\begin{frame}{Silhouette Coefficient}
\textbf{Silhouette Coefficient for observation i:}

\[s(i) = \frac{b(i) - a(i)}{\max(a(i),b(i))}\]

where:

\begin{itemize}
\item
  \(a(i)\) = avg distance to points in same cluster
\item
  \(b(i)\) = avg distance to points in nearest neighboring cluster
\end{itemize}
\end{frame}

\begin{frame}{Silhouette Interpretation}
\begin{itemize}
\item
  \(s(i) \approx + 1\): Well-matched to cluster
\item
  \(s(i) \approx 0\): On border between clusters
\item
  \(s(i) \approx - 1\): Likely in wrong cluster
\end{itemize}
\end{frame}

\begin{frame}{Using Silhouette for Optimal k}
\textbf{Average Silhouette Width:}

\[\overline{s} = \frac{1}{n}\sum_{i = 1}^{n}s(i)\]
\end{frame}

\begin{frame}{Silhouette Procedure}
\begin{enumerate}
\item
  Run clustering for different k values
\item
  Calculate average silhouette width for each k
\item
  Choose k that maximizes \(\overline{s}\)
\end{enumerate}

\textbf{Advantage:} Provides both quality measure and optimal k
\end{frame}

\begin{frame}{Case Study: Silhouette Plot Insights}
\textbf{Silhouette Plot Reveals:}

\begin{itemize}
\item
  Most customers have positive coefficients (well-assigned)
\item
  All clusters extend beyond average (0.458)
\item
  Some variation in cluster quality
\item
  Few customers near zero (boundary cases)
\item
  Very few negative coefficients (rare misclassifications)
\end{itemize}

\textbf{Conclusion:} Generally good cluster quality across all 4
segments
\end{frame}

\begin{frame}{Case Study: Cluster Profiling Results}
\textbf{Four Distinct Customer Segments Identified:}

\begin{enumerate}
\item
  \textbf{Engaged Selective (15.3\%)}: High basket, high returns, low
  browsing
\item
  \textbf{Low-Value Browsers (35.4\%)}: Browse but don't buy much
\item
  \textbf{Premium High-Value (21.6\%)}: Highest spenders, frequent
  purchases
\item
  \textbf{Frequent Small-Basket (27.7\%)}: Regular small orders
\end{enumerate}

Each differs by \textgreater10\% from average on multiple dimensions
\end{frame}

\begin{frame}{Cluster Validation}
\textbf{Internal Validation} (using data only):

\begin{itemize}
\item
  Within-Cluster Sum of Squares (WCSS) - lower is better
\item
  Silhouette Coefficient - higher is better
\item
  Davies-Bouldin Index - lower is better
\item
  Dunn Index - higher is better
\end{itemize}
\end{frame}

\begin{frame}{Cluster Validation}
\textbf{External Validation} (when true labels available):

\begin{itemize}
\item
  Adjusted Rand Index (ARI)
\item
  Normalized Mutual Information (NMI)
\end{itemize}
\end{frame}

\begin{frame}{Davies-Bouldin Index}
Measures ratio of within-cluster dispersion to between-cluster
separation

\[\text{ DB } = \frac{1}{k}\sum_{i = 1}^{k}\max\limits_{j \neq i}\left( \frac{\sigma_{i} + \sigma_{j}}{d\left( c_{i},c_{j} \right)} \right)\]
\end{frame}

\begin{frame}{Davies-Bouldin Index}
\textbf{Interpretation:}

\begin{itemize}
\item
  Lower values indicate better clustering
\item
  Compact clusters that are far apart
\item
  Can compare different k values or methods
\end{itemize}
\end{frame}

\begin{frame}{Curse of Dimensionality}
As dimensions (p) increase, problems arise
\end{frame}

\begin{frame}{Curse of Dimensionality: Problems}
\begin{enumerate}
\item
  Distance becomes less meaningful (all points appear equidistant)
\item
  Data becomes sparse (observations spread out)
\item
  Computational cost increases dramatically
\end{enumerate}
\end{frame}

\begin{frame}{Curse of Dimensionality: Solutions}
\begin{itemize}
\item
  Use PCA or feature selection before clustering
\item
  Select only relevant variables
\item
  Use specialized high-dimensional algorithms
\end{itemize}

\textbf{Rule:} If p is large relative to n, reduce dimensions first
\end{frame}

\begin{frame}{Handling Outliers}
\textbf{Impact by Method:}

{\def\LTcaptype{none} % do not increment counter
\begin{longtable}[]{@{}
  >{\raggedright\arraybackslash}p{(\linewidth - 2\tabcolsep) * \real{0.6000}}
  >{\centering\arraybackslash}p{(\linewidth - 2\tabcolsep) * \real{0.4000}}@{}}
\toprule\noalign{}
\endhead
\textbf{Method} & \textbf{Sensitivity} \\
K-means & High \\
Ward's Method & High \\
Single Linkage & Medium \\
K-medoids & Low (Robust) \\
\bottomrule\noalign{}
\end{longtable}
}
\end{frame}

\begin{frame}{Handling Outliers: Strategies}
\begin{itemize}
\item
  Pre-processing: Detect and remove outliers
\item
  Use robust methods (k-medoids)
\item
  Accept outlier clusters
\end{itemize}
\end{frame}

\begin{frame}{When to Use Hierarchical Clustering}
\begin{itemize}
\item
  Small to medium datasets (n \textless{} 5,000)
\item
  Want to explore different k values
\item
  Need hierarchical structure
\item
  Don't know k in advance
\end{itemize}
\end{frame}

\begin{frame}{When to Use K-Means}
\begin{itemize}
\item
  Large datasets (n \textgreater{} 5,000)
\item
  Approximately know k
\item
  Need speed and efficiency
\item
  Clusters roughly spherical
\end{itemize}
\end{frame}

\begin{frame}{Cluster Analysis Workflow}
\begin{enumerate}
\item
  \textbf{Define objective} - What questions to answer?
\item
  \textbf{Select variables} - Domain knowledge
\item
  \textbf{Preprocess data} - Handle missing values, outliers
\item
  \textbf{Standardize} - If variables on different scales
\item
  \textbf{Choose method} - Based on data characteristics
\end{enumerate}
\end{frame}

\begin{frame}{Cluster Analysis Workflow}
\begin{enumerate}
\setcounter{enumi}{5}
\item
  \textbf{Determine k} - Multiple criteria
\item
  \textbf{Run clustering} - Multiple times for k-means
\item
  \textbf{Validate results} - Internal and stability checks
\item
  \textbf{Interpret clusters} - Profile and name clusters
\item
  \textbf{Refine and iterate} - Based on insights
\end{enumerate}
\end{frame}

\begin{frame}{Common Pitfalls to Avoid}
\begin{enumerate}
\item
  \textbf{Not standardizing} when variables have different scales
\item
  \textbf{Using k-means} with non-spherical clusters
\item
  \textbf{Ignoring outliers} - can severely distort results
\item
  \textbf{Over-interpreting} - clustering always finds structure, even
  in random data
\end{enumerate}
\end{frame}

\begin{frame}{Common Pitfalls to Avoid}
\begin{enumerate}
\setcounter{enumi}{4}
\item
  \textbf{Using too many variables} - curse of dimensionality
\item
  \textbf{Running k-means once} - try multiple initializations
\item
  \textbf{Choosing k without validation} - use multiple methods
\end{enumerate}
\end{frame}

\begin{frame}{Best Practices}
\begin{enumerate}
\item
  \textbf{Try multiple methods} - Compare hierarchical, k-means, etc.
\item
  \textbf{Validate stability} - Bootstrap samples, different
  initializations
\item
  \textbf{Visualize extensively} - Scatter plots, dendrograms, parallel
  coordinates
\end{enumerate}
\end{frame}

\begin{frame}{Best Practices}
\begin{enumerate}
\setcounter{enumi}{3}
\item
  \textbf{Use domain knowledge} - Statistical metrics + practical sense
\item
  \textbf{Document decisions} - Why certain methods, parameters chosen
\item
  \textbf{Check interpretability} - Can you explain and use clusters?
\end{enumerate}
\end{frame}

\begin{frame}{Key Takeaways: Fundamental Concepts}
\begin{itemize}
\item
  Cluster analysis discovers natural groupings (unsupervised)
\item
  Distance measures are crucial (Euclidean, Manhattan)
\item
  Standardization essential for different scales
\end{itemize}

\textbf{Case Study:} Standardizing customer data (dollars, counts,
proportions) ensured equal feature contribution
\end{frame}

\begin{frame}{Key Takeaways: Methods}
\begin{itemize}
\item
  Hierarchical: Creates tree structure, multiple k values
\item
  K-means: Fast, scalable, requires specifying k
\item
  K-medoids: Robust alternative to k-means
\end{itemize}

\textbf{Case Study:} Both hierarchical and k-means identified 4 customer
segments with identical silhouette scores (0.458)
\end{frame}

\begin{frame}{Key Takeaways: Validation}
\begin{itemize}
\item
  Elbow method and silhouette analysis for optimal k
\item
  Multiple validation measures for quality assessment
\end{itemize}

\textbf{Case Study:} Convergence of dendrogram, elbow, and silhouette
all suggested k=4 segments
\end{frame}

\begin{frame}{Summary: Method Selection Guide}
{\def\LTcaptype{none} % do not increment counter
\begin{longtable}[]{@{}
  >{\raggedright\arraybackslash}p{(\linewidth - 2\tabcolsep) * \real{0.3333}}
  >{\raggedright\arraybackslash}p{(\linewidth - 2\tabcolsep) * \real{0.6667}}@{}}
\toprule\noalign{}
\endhead
\textbf{Situation} & \textbf{Recommended Method} \\
Small dataset (n \textless{} 1,000) & Hierarchical (Ward's or
Average) \\
Large dataset (n \textgreater{} 10,000) & K-means with k-means++ \\
Outliers present & K-medoids or preprocessing \\
Non-spherical clusters & DBSCAN or hierarchical \\
\bottomrule\noalign{}
\end{longtable}
}
\end{frame}

\begin{frame}{Summary: Method Selection Guide}
{\def\LTcaptype{none} % do not increment counter
\begin{longtable}[]{@{}
  >{\raggedright\arraybackslash}p{(\linewidth - 2\tabcolsep) * \real{0.3333}}
  >{\raggedright\arraybackslash}p{(\linewidth - 2\tabcolsep) * \real{0.6667}}@{}}
\toprule\noalign{}
\endhead
\textbf{Situation} & \textbf{Recommended Method} \\
Don't know k & Hierarchical, then elbow/silhouette \\
High dimensions & PCA first, then k-means \\
Mixed data types & Gower distance with hierarchical \\
\bottomrule\noalign{}
\end{longtable}
}
\end{frame}

\begin{frame}{Advanced Topics (Beyond This Course)}
\textbf{Density-Based Methods:}

\begin{itemize}
\item
  DBSCAN - finds arbitrary shapes, identifies outliers
\end{itemize}

\textbf{Model-Based:}

\begin{itemize}
\item
  Gaussian Mixture Models (GMM) - probabilistic approach
\end{itemize}
\end{frame}

\begin{frame}{Advanced Topics (Beyond This Course)}
\textbf{Fuzzy Clustering:}

\begin{itemize}
\item
  Soft assignment (membership degrees)
\end{itemize}

\textbf{Subspace Clustering:}

\begin{itemize}
\item
  For high-dimensional data, different subspaces
\end{itemize}
\end{frame}

\begin{frame}{Real-World Applications: Business}
\textbf{Marketing \& Business:}

\begin{itemize}
\item
  Customer segmentation for targeted marketing
\item
  Product recommendation systems
\item
  Market basket analysis
\end{itemize}
\end{frame}

\begin{frame}{Real-World Applications: Healthcare}
\textbf{Healthcare:}

\begin{itemize}
\item
  Patient stratification for personalized medicine
\item
  Disease subtype identification
\item
  Medical image segmentation
\end{itemize}
\end{frame}

\begin{frame}{Real-World Applications: Finance}
\textbf{Finance:}

\begin{itemize}
\item
  Fraud detection and anomaly identification
\item
  Credit risk assessment
\item
  Portfolio diversification
\end{itemize}
\end{frame}

\begin{frame}{Example: Customer Segmentation}
\textbf{Scenario:} E-commerce company with 100,000 customers

\textbf{Variables:}

\begin{itemize}
\item
  Purchase frequency
\item
  Average order value
\item
  Product category preferences
\item
  Time since last purchase
\item
  Customer lifetime value
\end{itemize}
\end{frame}

\begin{frame}{Example: Customer Segmentation Process}
\begin{enumerate}
\item
  Standardize variables (different scales)
\item
  Try k-means for k = 2 to 10
\item
  Use elbow method and silhouette analysis
\item
  Identify k = 5 optimal clusters
\item
  Profile each segment
\item
  Develop targeted marketing strategies
\end{enumerate}
\end{frame}

\begin{frame}{Recommended Resources: Books}
\textbf{Textbooks:}

\begin{itemize}
\item
  Everitt et al. (2011) - Cluster Analysis (5th ed.)
\item
  James et al. (2021) - Introduction to Statistical Learning
\end{itemize}
\end{frame}

\begin{frame}{Recommended Resources: Software}
\textbf{Software:}

\begin{itemize}
\item
  Python: scikit-learn (KMeans, AgglomerativeClustering)
\item
  R: stats package (kmeans, hclust)
\end{itemize}
\end{frame}

\begin{frame}{Recommended Resources: Online Videos}
\textbf{StatQuest YouTube Channel:}

\begin{enumerate}
\item
  \textbf{K-Means Clustering:}\\
  \url{https://www.youtube.com/watch?v=4b5d3muPQmA}
\item
  \textbf{Hierarchical Clustering:}\\
  \url{https://www.youtube.com/watch?v=7xHsRkOdVwo}
\item
  \textbf{Validation Methods (Elbow \& Silhouette):}\\
  \url{https://www.youtube.com/watch?v=DzrvLpxTxJw}
\end{enumerate}
\end{frame}

\begin{frame}{Recommended Resources: Additional Online}
\textbf{Additional Online Resources:}

\begin{itemize}
\item
  Scikit-learn documentation
\item
  Coursera/edX courses on unsupervised learning
\item
  Interactive clustering visualizations
\end{itemize}
\end{frame}

\begin{frame}{Part 2: Practical Application}
\textbf{ Customer Segmentation Case Study }

E-Commerce Cluster Analysis Using Hierarchical and K-Means Methods
\end{frame}

\begin{frame}{Business Context}
\textbf{Business Problem:} An e-commerce company seeks to understand
their customer base by discovering natural segments based on purchasing
behavior, engagement patterns, and browsing habits.

\textbf{Key Questions:}

\begin{itemize}
\item
  How many distinct customer segments exist in our data?
\item
  What behavioral patterns characterize each segment?
\item
  How can we tailor marketing strategies to each discovered segment?
\end{itemize}
\end{frame}

\begin{frame}{Dataset Overview}
\textbf{Dataset:} 2,000 customers with 7 behavioral features

\textbf{Variables:}

\begin{itemize}
\item
  monthly\_purchases
\item
  avg\_basket\_size
\item
  total\_spend
\item
  session\_duration
\item
  email\_clicks
\item
  product\_views
\item
  return\_rate
\end{itemize}

\textbf{Important:} No predefined labels (unsupervised learning)
\end{frame}

\begin{frame}{Exploratory Data Analysis}
\begin{block}{Why EDA Matters}
Before attempting to discover customer segments, we must first
understand the characteristics and relationships within our data. Unlike
supervised learning where we have predefined labels, cluster analysis is
exploratory in nature, making this preliminary investigation even more
critical.
\end{block}
\end{frame}

\begin{frame}{EDA: Approach}
\textbf{What We Do:}

\begin{itemize}
\item
  Visualize distributions using histograms
\item
  Identify skewness, outliers, and typical value ranges
\item
  Compute correlation matrix to understand relationships
\item
  Assess variable redundancy
\end{itemize}

\textbf{Why:} Understanding data structure before algorithmic analysis
helps anticipate which features might drive segmentation.
\end{frame}

\begin{frame}[fragile]{EDA: Distribution Analysis}
\begin{Shaded}
\begin{Highlighting}[]
\CommentTok{\# Load customer data}
\NormalTok{df }\OperatorTok{=}\NormalTok{ pd.read\_csv(}\StringTok{"customer\_data.csv"}\NormalTok{)}

\CommentTok{\# Summary statistics}
\BuiltInTok{print}\NormalTok{(df.describe())}

\CommentTok{\# Distribution plots for each variable}
\NormalTok{fig, axes }\OperatorTok{=}\NormalTok{ plt.subplots(}\DecValTok{3}\NormalTok{, }\DecValTok{3}\NormalTok{, figsize}\OperatorTok{=}\NormalTok{(}\DecValTok{15}\NormalTok{, }\DecValTok{12}\NormalTok{))}
\ControlFlowTok{for}\NormalTok{ idx, col }\KeywordTok{in} \BuiltInTok{enumerate}\NormalTok{(df.columns):}
\NormalTok{    axes[row, col\_idx].hist(df[col], bins}\OperatorTok{=}\DecValTok{30}\NormalTok{)}
\end{Highlighting}
\end{Shaded}
\end{frame}

\begin{frame}[fragile]{EDA: Correlation Analysis}
\begin{Shaded}
\begin{Highlighting}[]
\CommentTok{\# Correlation matrix}
\NormalTok{correlation\_matrix }\OperatorTok{=}\NormalTok{ df.corr()}
\NormalTok{sns.heatmap(correlation\_matrix, annot}\OperatorTok{=}\VariableTok{True}\NormalTok{,}
\NormalTok{            fmt}\OperatorTok{=}\StringTok{\textquotesingle{}.2f\textquotesingle{}}\NormalTok{, cmap}\OperatorTok{=}\StringTok{\textquotesingle{}coolwarm\textquotesingle{}}\NormalTok{)}

\CommentTok{\# Identify strongest correlations}
\CommentTok{\# Example: monthly\_purchases \textless{}{-}\textgreater{} total\_spend: 0.94}
\CommentTok{\#          avg\_basket\_size \textless{}{-}\textgreater{} total\_spend: 0.89}
\end{Highlighting}
\end{Shaded}
\end{frame}

\begin{frame}{EDA: Key Findings}
\textbf{Outcome:}

\begin{itemize}
\item
  Balanced dataset with diverse customer behaviors
\item
  Right-skewed distributions (some extreme high spenders)
\item
  Strong positive correlations between purchase-related variables
\item
  Email clicks correlate moderately with purchase frequency
\end{itemize}

\textbf{Implication:} Multiple distinct behavioral patterns suggest
natural customer segments exist.
\end{frame}

\begin{frame}{Data Standardization}
\begin{block}{Why Standardization is Critical}
Cluster analysis algorithms rely on distance metrics to measure
similarity. However, our behavioral variables are measured in different
units and scales:

\begin{itemize}
\item
  Purchases (counts)
\item
  Spending (dollars: 58 to 7,891)
\item
  Return rate (proportions: 0.0 to 0.5)
\end{itemize}
\end{block}
\end{frame}

\begin{frame}{The Problem Without Standardization}
\textbf{Without standardization:} Variables with larger numeric ranges
dominate distance calculations.

\textbf{Example:}

\begin{itemize}
\item
  Customer A: 1 purchase, 100 dollars
\item
  Customer B: 2 purchases, 200 dollars
\end{itemize}

Distance dominated by dollar difference (100) rather than purchase
difference (1).

This leads to biased clustering that reflects scale differences, not
true behavioral patterns.
\end{frame}

\begin{frame}{Standardization: Approach}
\textbf{Z-score Standardization:}

\[z_{i} = \frac{x_{i} - \mu}{\sigma}\]

\textbf{Properties:}

\begin{itemize}
\item
  Transforms to mean = 0, standard deviation = 1
\item
  Preserves distribution shape
\item
  Ensures equal contribution to distance calculations
\item
  Value of 2.0 means ``2 std deviations above mean''
\end{itemize}
\end{frame}

\begin{frame}[fragile]{Standardization: Implementation}
\begin{Shaded}
\begin{Highlighting}[]
\ImportTok{from}\NormalTok{ sklearn.preprocessing }\ImportTok{import}\NormalTok{ StandardScaler}

\CommentTok{\# Standardize features}
\NormalTok{scaler }\OperatorTok{=}\NormalTok{ StandardScaler()}
\NormalTok{X\_standardized }\OperatorTok{=}\NormalTok{ scaler.fit\_transform(df)}

\CommentTok{\# Convert back to DataFrame}
\NormalTok{df\_standardized }\OperatorTok{=}\NormalTok{ pd.DataFrame(}
\NormalTok{    X\_standardized,}
\NormalTok{    columns}\OperatorTok{=}\NormalTok{df.columns}
\NormalTok{)}

\BuiltInTok{print}\NormalTok{(df\_standardized.describe())}
\CommentTok{\# All means approx 0, all std approx 1}
\end{Highlighting}
\end{Shaded}
\end{frame}

\begin{frame}{Standardization: Outcome}
\textbf{Result:} All variables successfully transformed to mean
approximately 0 and standard deviation approximately 1.

\textbf{Impact:} Clustering algorithms now treat all behavioral
dimensions equally when computing distances between customers.

\textbf{Example:} Total spend (originally 58-7,891 dollars) and return
rate (originally 0.0-0.5) now contribute equally to customer similarity.
\end{frame}

\begin{frame}{Hierarchical Clustering}
\begin{block}{The Fundamental Question}
\textbf{How many customer segments should we look for?}

Unlike supervised classification where the number of classes is
predetermined, clustering requires us to discover the appropriate number
of groups.

\textbf{Hierarchical clustering provides an elegant solution:} Build a
complete hierarchy of nested clusters without specifying k in advance.
\end{block}
\end{frame}

\begin{frame}{Hierarchical Clustering: Advantage}
\textbf{Key Benefit:} Creates a tree-like structure (dendrogram) showing
how customers progressively merge into larger groups.

\textbf{Result:} We can identify the most natural number of segments
based on where large jumps in distance occur.

No need to predefine the number of clusters
\end{frame}

\begin{frame}{Hierarchical Clustering: Linkage Methods}
\textbf{How do we measure distance between clusters?}

\begin{itemize}
\item
  \textbf{Single Linkage:} Minimum distance between any two points
\item
  \textbf{Complete Linkage:} Maximum distance between any two points
\item
  \textbf{Average Linkage:} Mean distance between all pairs
\item
  \textbf{Ward's Method:} Minimizes within-cluster variance
\end{itemize}

\textbf{Recommendation:} Ward's method typically produces the most
balanced clusters for customer segmentation.
\end{frame}

\begin{frame}[fragile]{Hierarchical Clustering: Implementation}
\begin{Shaded}
\begin{Highlighting}[]
\ImportTok{from}\NormalTok{ scipy.cluster.hierarchy }\ImportTok{import}\NormalTok{ linkage, dendrogram}

\CommentTok{\# Compute linkage matrices for different methods}
\NormalTok{linkage\_methods }\OperatorTok{=}\NormalTok{ [}\StringTok{\textquotesingle{}single\textquotesingle{}}\NormalTok{, }\StringTok{\textquotesingle{}complete\textquotesingle{}}\NormalTok{,}
                   \StringTok{\textquotesingle{}average\textquotesingle{}}\NormalTok{, }\StringTok{\textquotesingle{}ward\textquotesingle{}}\NormalTok{]}
\NormalTok{linkage\_matrices }\OperatorTok{=}\NormalTok{ \{\}}

\ControlFlowTok{for}\NormalTok{ method }\KeywordTok{in}\NormalTok{ linkage\_methods:}
\NormalTok{    linkage\_matrices[method] }\OperatorTok{=}\NormalTok{ linkage(}
\NormalTok{        X\_standardized,}
\NormalTok{        method}\OperatorTok{=}\NormalTok{method}
\NormalTok{    )}
\end{Highlighting}
\end{Shaded}
\end{frame}

\begin{frame}[fragile]{Dendrogram Interpretation}
\begin{Shaded}
\begin{Highlighting}[]
\CommentTok{\# Create dendrogram for Ward\textquotesingle{}s method}
\NormalTok{plt.figure(figsize}\OperatorTok{=}\NormalTok{(}\DecValTok{14}\NormalTok{, }\DecValTok{7}\NormalTok{))}
\NormalTok{dendrogram(linkage\_matrices[}\StringTok{\textquotesingle{}ward\textquotesingle{}}\NormalTok{],}
\NormalTok{           no\_labels}\OperatorTok{=}\VariableTok{True}\NormalTok{,}
\NormalTok{           color\_threshold}\OperatorTok{=}\DecValTok{50}\NormalTok{)}
\NormalTok{plt.axhline(y}\OperatorTok{=}\DecValTok{50}\NormalTok{, color}\OperatorTok{=}\StringTok{\textquotesingle{}r\textquotesingle{}}\NormalTok{, linestyle}\OperatorTok{=}\StringTok{\textquotesingle{}{-}{-}\textquotesingle{}}\NormalTok{,}
\NormalTok{            label}\OperatorTok{=}\StringTok{\textquotesingle{}Potential cut (4 clusters)\textquotesingle{}}\NormalTok{)}
\end{Highlighting}
\end{Shaded}

\textbf{How to Read:}

\begin{itemize}
\item
  Horizontal axis: Customers
\item
  Vertical axis: Distance at which clusters merge
\item
  Large vertical gaps suggest natural cluster boundaries
\end{itemize}
\end{frame}

\begin{frame}{Hierarchical Clustering: Outcome}
\textbf{Finding:} Four dendrograms reveal distinct clustering structures
depending on linkage method.

\textbf{Observations:}

\begin{itemize}
\item
  Single linkage: Chaining effects with unbalanced clusters
\item
  Complete linkage: Too many small clusters
\item
  Average linkage: Compromise between extremes
\item
  Ward's method: Clear hierarchical structure, balanced sizes
\end{itemize}

\textbf{Decision:} Ward's dendrogram suggests 4 clusters as a natural
choice.
\end{frame}

\begin{frame}[fragile,shrink]{Extracting Clusters from Dendrogram}
\begin{Shaded}
\begin{Highlighting}[]
\ImportTok{from}\NormalTok{ scipy.cluster.hierarchy }\ImportTok{import}\NormalTok{ fcluster}

\CommentTok{\# Extract 4 clusters using Ward\textquotesingle{}s method}
\NormalTok{n\_clusters\_hier }\OperatorTok{=} \DecValTok{4}
\NormalTok{hierarchical\_labels }\OperatorTok{=}\NormalTok{ fcluster(}
\NormalTok{    linkage\_matrices[}\StringTok{\textquotesingle{}ward\textquotesingle{}}\NormalTok{],}
\NormalTok{    n\_clusters\_hier,}
\NormalTok{    criterion}\OperatorTok{=}\StringTok{\textquotesingle{}maxclust\textquotesingle{}}
\NormalTok{)}

\CommentTok{\# Calculate silhouette score}
\NormalTok{silhouette\_hier }\OperatorTok{=}\NormalTok{ silhouette\_score(}
\NormalTok{    X\_standardized,}
\NormalTok{    hierarchical\_labels}
\NormalTok{)}
\CommentTok{\# Result: 0.458 (moderate cluster quality)}
\end{Highlighting}
\end{Shaded}
\end{frame}

\begin{frame}{Hierarchical Clustering: Results}
\textbf{Cluster Sizes:}

\begin{itemize}
\item
  Cluster 0: 440 customers (22.0\%)
\item
  Cluster 1: 300 customers (15.0\%)
\item
  Cluster 2: 560 customers (28.0\%)
\item
  Cluster 3: 700 customers (35.0\%)
\end{itemize}

\textbf{Silhouette Score:} 0.458

\textbf{Interpretation:} Customers reasonably well-separated into
clusters with balanced sizes. Score above 0.4 is acceptable for customer
segmentation where behavioral boundaries are often fuzzy.
\end{frame}

\begin{frame}{K-Means Clustering}
\begin{block}{Why K-Means?}
While hierarchical clustering provided valuable insights through
dendrogram visualization, it has computational limitations:

\begin{itemize}
\item
  Complexity: O(n squared) or worse
\item
  Impractical for very large datasets
\end{itemize}

\textbf{K-means offers a scalable alternative:}

\begin{itemize}
\item
  Works well with larger customer bases
\item
  Computational complexity: O(n times k times p times iterations)
\end{itemize}
\end{block}
\end{frame}

\begin{frame}{K-Means: The Challenge}
\textbf{Requirement:} Must specify the number of clusters (k) in
advance.

\textbf{Solution:} Elbow method provides a data-driven approach.

\textbf{Goal:} Identify the point where adding more clusters provides
diminishing returns in terms of improved fit.
\end{frame}

\begin{frame}{Elbow Method: Approach}
K-means minimizes within-cluster sum of squared distances (inertia):

\[\text{ WCSS } = \sum_{i = 1}^{k}\sum_{x \in C_{i}}\| x - \mu_{i}\|^{2}\]

\textbf{Procedure:}

\begin{itemize}
\item
  Run clustering for k = 2 to 10
\item
  Calculate inertia for each k
\item
  Plot inertia vs. k
\item
  Look for the ``elbow'' point
\end{itemize}
\end{frame}

\begin{frame}[fragile]{Elbow Method: Implementation}
\begin{Shaded}
\begin{Highlighting}[]
\ImportTok{from}\NormalTok{ sklearn.cluster }\ImportTok{import}\NormalTok{ KMeans}

\CommentTok{\# Test different k values}
\NormalTok{inertias }\OperatorTok{=}\NormalTok{ []}
\NormalTok{silhouette\_scores }\OperatorTok{=}\NormalTok{ []}
\NormalTok{K\_range }\OperatorTok{=} \BuiltInTok{range}\NormalTok{(}\DecValTok{2}\NormalTok{, }\DecValTok{11}\NormalTok{)}

\ControlFlowTok{for}\NormalTok{ k }\KeywordTok{in}\NormalTok{ K\_range:}
\NormalTok{    kmeans }\OperatorTok{=}\NormalTok{ KMeans(n\_clusters}\OperatorTok{=}\NormalTok{k, random\_state}\OperatorTok{=}\DecValTok{42}\NormalTok{)}
\NormalTok{    kmeans.fit(X\_standardized)}
\NormalTok{    inertias.append(kmeans.inertia\_)}
\NormalTok{    silhouette\_scores.append(}
\NormalTok{        silhouette\_score(X\_standardized, kmeans.labels\_)}
\NormalTok{    )}
\end{Highlighting}
\end{Shaded}
\end{frame}

\begin{frame}{Elbow Method: Results}
\textbf{Plot Analysis:}

\begin{itemize}
\item
  Elbow curve shows diminishing returns after k=4
\item
  Silhouette scores peak at k=4
\item
  Both metrics converge on k=4
\end{itemize}

\textbf{Decision:} Both hierarchical and elbow analyses suggest k=4
\end{frame}

\begin{frame}[fragile]{K-Means: Final Clustering}
\begin{Shaded}
\begin{Highlighting}[]
\CommentTok{\# Apply k{-}means with optimal k}
\NormalTok{optimal\_k }\OperatorTok{=} \DecValTok{4}
\NormalTok{kmeans\_final }\OperatorTok{=}\NormalTok{ KMeans(}
\NormalTok{    n\_clusters}\OperatorTok{=}\NormalTok{optimal\_k,}
\NormalTok{    random\_state}\OperatorTok{=}\DecValTok{42}\NormalTok{,}
\NormalTok{    n\_init}\OperatorTok{=}\DecValTok{10}
\NormalTok{)}
\NormalTok{kmeans\_labels }\OperatorTok{=}\NormalTok{ kmeans\_final.fit\_predict(}
\NormalTok{    X\_standardized}
\NormalTok{)}

\CommentTok{\# Silhouette score: 0.458 (same as hierarchical)}
\end{Highlighting}
\end{Shaded}
\end{frame}

\begin{frame}{K-Means: Convergence Validation}
\textbf{Key Finding:} K-means achieves identical silhouette score to
hierarchical clustering (0.458).

\textbf{Significance:} Convergence between two fundamentally different
algorithms provides strong evidence that four customer segments
represent genuine structure in the data.

\textbf{Confidence:} Clusters reflect real behavioral patterns rather
than algorithmic artifacts.

\textbf{Recommendation:} Use k-means for production deployment due to
computational efficiency.
\end{frame}

\begin{frame}{Cluster Interpretation and Profiling}
\begin{block}{From Statistics to Business Value}
Successfully identifying clusters is only the first step. The real
business value comes from understanding what distinguishes each segment
and translating these differences into actionable marketing strategies.

\textbf{Goal:} Transform abstract cluster labels into concrete customer
personas.
\end{block}
\end{frame}

\begin{frame}[fragile]{Interpretation: Approach}
\begin{Shaded}
\begin{Highlighting}[]
\CommentTok{\# Add cluster labels to original data}
\NormalTok{df\_with\_clusters }\OperatorTok{=}\NormalTok{ df.copy()}
\NormalTok{df\_with\_clusters[}\StringTok{\textquotesingle{}Cluster\_KMeans\textquotesingle{}}\NormalTok{] }\OperatorTok{=}\NormalTok{ kmeans\_labels}

\CommentTok{\# Calculate cluster means (original units)}
\NormalTok{cluster\_profiles }\OperatorTok{=}\NormalTok{ df\_with\_clusters.groupby(}
    \StringTok{\textquotesingle{}Cluster\_KMeans\textquotesingle{}}
\NormalTok{)[df.columns].mean()}

\CommentTok{\# Create heatmap}
\NormalTok{sns.heatmap(cluster\_profiles.T, annot}\OperatorTok{=}\VariableTok{True}\NormalTok{)}
\end{Highlighting}
\end{Shaded}

\textbf{Why original units?} Makes profiles interpretable to business
stakeholders.
\end{frame}

\begin{frame}[fragile]{Cluster Characterization: Method}
\begin{Shaded}
\begin{Highlighting}[]
\CommentTok{\# Compare each cluster to overall means}
\NormalTok{overall\_means }\OperatorTok{=}\NormalTok{ df.mean()}

\ControlFlowTok{for}\NormalTok{ cluster\_id }\KeywordTok{in} \BuiltInTok{range}\NormalTok{(optimal\_k):}
\NormalTok{    cluster\_mean }\OperatorTok{=}\NormalTok{ cluster\_profiles.loc[cluster\_id]}

    \CommentTok{\# Calculate percentage differences}
\NormalTok{    differences }\OperatorTok{=}\NormalTok{ (}
\NormalTok{        (cluster\_mean }\OperatorTok{{-}}\NormalTok{ overall\_means) }\OperatorTok{/}
\NormalTok{        overall\_means }\OperatorTok{*} \DecValTok{100}
\NormalTok{    )}

    \CommentTok{\# Identify distinctive features (\textgreater{}10\% difference)}
\NormalTok{    high\_features }\OperatorTok{=}\NormalTok{ differences.nlargest(}\DecValTok{3}\NormalTok{)}
\NormalTok{    low\_features }\OperatorTok{=}\NormalTok{ differences.nsmallest(}\DecValTok{3}\NormalTok{)}
\end{Highlighting}
\end{Shaded}
\end{frame}

\begin{frame}{Cluster 0: Engaged but Selective Shoppers}
\textbf{Size:} 307 customers (15.3\%)

\textbf{Distinctive High Features:}

\begin{itemize}
\item
  avg\_basket\_size: +72.0\% vs average
\item
  return\_rate: +56.2\% vs average
\item
  email\_clicks: +51.9\% vs average
\end{itemize}

\textbf{Distinctive Low Features:}

\begin{itemize}
\item
  session\_duration: -69.6\% vs average
\item
  product\_views: -58.0\% vs average
\item
  monthly\_purchases: -42.1\% vs average
\end{itemize}
\end{frame}

\begin{frame}{Cluster 1: Low-Value Browsers}
\textbf{Size:} 707 customers (35.4\%)

\textbf{Distinctive High Features:}

\begin{itemize}
\item
  session\_duration: +19.0\% vs average
\item
  return\_rate: +15.2\% vs average
\end{itemize}

\textbf{Distinctive Low Features:}

\begin{itemize}
\item
  total\_spend: -83.8\% vs average
\item
  email\_clicks: -78.3\% vs average
\item
  monthly\_purchases: -76.2\% vs average
\end{itemize}

\textbf{Characterization:} Spend time browsing but make few purchases.
\end{frame}

\begin{frame}{Cluster 2: Premium High-Value Customers}
\textbf{Size:} 432 customers (21.6\%)

\textbf{Distinctive High Features:}

\begin{itemize}
\item
  total\_spend: +168.3\% vs average
\item
  avg\_basket\_size: +127.5\% vs average
\item
  monthly\_purchases: +124.1\% vs average
\end{itemize}

\textbf{Distinctive Low Features:}

\begin{itemize}
\item
  return\_rate: -24.1\% vs average
\end{itemize}

\textbf{Characterization:} The premium segment with highest spending and
purchase frequency.
\end{frame}

\begin{frame}{Cluster 3: Frequent Small-Basket Shoppers}
\textbf{Size:} 554 customers (27.7\%)

\textbf{Distinctive High Features:}

\begin{itemize}
\item
  monthly\_purchases: +23.8\% vs average
\item
  product\_views: +12.4\% vs average
\end{itemize}

\textbf{Distinctive Low Features:}

\begin{itemize}
\item
  avg\_basket\_size: -48.8\% vs average
\item
  total\_spend: -45.1\% vs average
\item
  return\_rate: -31.8\% vs average
\end{itemize}

\textbf{Characterization:} Regular shoppers with lower average order
values.
\end{frame}

\begin{frame}{Cluster Profiling: Outcome}
\textbf{Four Distinct Segments Identified:}

Each segment has distinctive characteristics differing by more than 10\%
from overall average on multiple dimensions.

\textbf{Confirmation:} These represent meaningfully different customer
types requiring differentiated marketing strategies.

\textbf{Next Step:} Develop targeted campaigns for each segment.
\end{frame}

\begin{frame}{Cluster Validation with Silhouette Analysis}
\begin{block}{Beyond Average Scores}
While we have identified interpretable customer segments, we should
validate the quality of our clustering solution.

\textbf{Questions:}

\begin{itemize}
\item
  Are customers well-matched to their assigned clusters?
\item
  Are some clusters poorly defined?
\item
  Are there misclassifications?
\end{itemize}
\end{block}
\end{frame}

\begin{frame}{Silhouette Analysis: Approach}
\textbf{Silhouette Coefficient for each customer:}

\[s(i) = \frac{b(i) - a(i)}{\max(a(i),b(i))}\]

where:

\begin{itemize}
\item
  \(a(i)\) = average distance to points in same cluster
\item
  \(b(i)\) = average distance to points in nearest neighboring cluster
\end{itemize}

\textbf{Range:} -1 to +1

\begin{itemize}
\item
  Near +1: Well-matched to cluster
\item
  Near 0: On border between clusters
\item
  Negative: Possible misclassification
\end{itemize}
\end{frame}

\begin{frame}[fragile]{Silhouette Plot: Implementation}
\begin{Shaded}
\begin{Highlighting}[]
\ImportTok{from}\NormalTok{ sklearn.metrics }\ImportTok{import}\NormalTok{ silhouette\_samples}

\CommentTok{\# Calculate silhouette values for each customer}
\NormalTok{silhouette\_vals }\OperatorTok{=}\NormalTok{ silhouette\_samples(}
\NormalTok{    X\_standardized,}
\NormalTok{    kmeans\_labels}
\NormalTok{)}

\CommentTok{\# Create visualization}
\ControlFlowTok{for}\NormalTok{ i }\KeywordTok{in} \BuiltInTok{range}\NormalTok{(optimal\_k):}
\NormalTok{    cluster\_vals }\OperatorTok{=}\NormalTok{ silhouette\_vals[kmeans\_labels }\OperatorTok{==}\NormalTok{ i]}
\NormalTok{    cluster\_vals.sort()}

\NormalTok{    plt.fill\_betweenx(y\_range, }\DecValTok{0}\NormalTok{, cluster\_vals)}
\end{Highlighting}
\end{Shaded}
\end{frame}

\begin{frame}{Silhouette Analysis: Interpretation}
\textbf{Key Observations:}

\begin{itemize}
\item
  Width of each section represents cluster size
\item
  All clusters extend beyond average score (0.458)
\item
  Most customers have positive coefficients
\item
  Few customers near zero (boundary cases)
\item
  Very few negative coefficients (rare misclassifications)
\end{itemize}

\textbf{Conclusion:} Generally good cluster quality across all four
segments. Segmentation is sound.
\end{frame}

\begin{frame}{Visualization in 2D Space}
\begin{block}{The Challenge}
We have worked with seven-dimensional customer data, which is impossible
to visualize directly.

\textbf{Problem:} Humans can only perceive 2-3 spatial dimensions
effectively.

\textbf{Solution:} Principal Component Analysis (PCA) projects the 7D
space onto 2D while preserving as much information as possible.
\end{block}
\end{frame}

\begin{frame}{PCA: Approach}
\textbf{What PCA Does:}

\begin{itemize}
\item
  Transforms 7 original variables into uncorrelated components
\item
  Orders components by variance explained
\item
  PC1 captures maximum variance
\item
  PC2 captures maximum remaining variance orthogonal to PC1
\end{itemize}

\textbf{Result:} 2D visualization retaining most important structure.
\end{frame}

\begin{frame}[fragile]{PCA Projection: Implementation}
\begin{Shaded}
\begin{Highlighting}[]
\ImportTok{from}\NormalTok{ sklearn.decomposition }\ImportTok{import}\NormalTok{ PCA}

\CommentTok{\# Apply PCA}
\NormalTok{pca }\OperatorTok{=}\NormalTok{ PCA(n\_components}\OperatorTok{=}\DecValTok{2}\NormalTok{)}
\NormalTok{X\_pca }\OperatorTok{=}\NormalTok{ pca.fit\_transform(X\_standardized)}

\CommentTok{\# Variance explained}
\BuiltInTok{print}\NormalTok{(}\SpecialStringTok{f"PC1: }\SpecialCharTok{\{}\NormalTok{pca}\SpecialCharTok{.}\NormalTok{explained\_variance\_ratio\_[}\DecValTok{0}\NormalTok{]}\SpecialCharTok{:.3f\}}\SpecialStringTok{"}\NormalTok{)}
\BuiltInTok{print}\NormalTok{(}\SpecialStringTok{f"PC2: }\SpecialCharTok{\{}\NormalTok{pca}\SpecialCharTok{.}\NormalTok{explained\_variance\_ratio\_[}\DecValTok{1}\NormalTok{]}\SpecialCharTok{:.3f\}}\SpecialStringTok{"}\NormalTok{)}

\CommentTok{\# Create scatter plot}
\NormalTok{plt.scatter(X\_pca[:, }\DecValTok{0}\NormalTok{], X\_pca[:, }\DecValTok{1}\NormalTok{],}
\NormalTok{            c}\OperatorTok{=}\NormalTok{kmeans\_labels, cmap}\OperatorTok{=}\StringTok{\textquotesingle{}tab10\textquotesingle{}}\NormalTok{)}
\end{Highlighting}
\end{Shaded}
\end{frame}

\begin{frame}{PCA Visualization: Outcome}
\textbf{Results:}

\begin{itemize}
\item
  2D projection captures moderate proportion of variance (40-60\%)
\item
  Reasonably good cluster separation visible
\item
  Some overlap consistent with silhouette analysis
\item
  K-means centroids clearly separated
\item
  Both methods show similar spatial patterns
\end{itemize}

\textbf{Interpretation:} While 2D view sacrifices some information, it
confirms clusters are spatially distinct and well-positioned.
\end{frame}

\begin{frame}{Business Recommendations}
\begin{block}{From Analysis to Action}
Having completed technical analysis and validation, we now translate
statistical findings into actionable business strategies.

\textbf{Goal:} Create customer personas and design specific marketing
tactics for each segment aligned with their behaviors and needs.
\end{block}
\end{frame}

\begin{frame}{Segment 1: Engaged but Selective Shoppers}
\textbf{Profile:} High basket sizes, high returns, engaged with emails
but low browsing.

\textbf{Marketing Strategy:}

\begin{itemize}
\item
  Launch VIP loyalty program with exclusive perks
\item
  Offer early access to new products
\item
  Provide free expedited shipping
\item
  Personalized product recommendations
\item
  Improve product information to reduce returns
\end{itemize}
\end{frame}

\begin{frame}{Segment 2: Low-Value Browsers}
\textbf{Profile:} High browsing time but low purchases and engagement.

\textbf{Marketing Strategy:}

\begin{itemize}
\item
  Implement abandoned cart recovery campaigns
\item
  Use retargeting ads to re-engage visitors
\item
  Offer limited-time discounts to incentivize first purchase
\item
  Improve product information and customer reviews
\item
  Create urgency with flash sales
\end{itemize}
\end{frame}

\begin{frame}{Segment 3: Premium High-Value Customers}
\textbf{Profile:} Highest spending, purchase frequency, and basket
sizes.

\textbf{Marketing Strategy:}

\begin{itemize}
\item
  Exclusive VIP treatment and recognition
\item
  Premium customer service channel
\item
  Early access to new collections
\item
  Referral program incentives
\item
  Maintain satisfaction to prevent churn
\end{itemize}
\end{frame}

\begin{frame}{Segment 4: Frequent Small-Basket Shoppers}
\textbf{Profile:} Regular purchases with lower average order values.

\textbf{Marketing Strategy:}

\begin{itemize}
\item
  Send targeted discount codes and bundle offers
\item
  Promote free shipping thresholds to increase basket size
\item
  Highlight clearance and sale items
\item
  Create value packs and multi-buy promotions
\item
  Build purchase frequency through regular engagement
\end{itemize}
\end{frame}

\begin{frame}{Key Findings: Summary}
\textbf{Methodology:}

\begin{itemize}
\item
  Analyzed 2,000 customers across 7 behavioral variables
\item
  Standardized data to ensure equal feature weighting
\item
  Applied both hierarchical (Ward's) and k-means clustering
\item
  Used elbow method and silhouette analysis
\end{itemize}

\textbf{Result:} Both methods converged on 4 distinct customer segments
with silhouette score of 0.458.
\end{frame}

\begin{frame}{Key Findings: Segments}
\textbf{Four Distinct Customer Segments:}

\begin{enumerate}
\item
  \textbf{Engaged Selective (15\%)}: High basket, high return
\item
  \textbf{Low-Value Browsers (35\%)}: Browse but don't buy
\item
  \textbf{Premium High-Value (22\%)}: Highest spenders
\item
  \textbf{Frequent Small-Basket (28\%)}: Regular small orders
\end{enumerate}

Each requires tailored marketing strategies.
\end{frame}

\begin{frame}{Methodological Learnings}
\textbf{Hierarchical Clustering:}

\begin{itemize}
\item
  Provides hierarchy view
\item
  No need to predefine k
\item
  Best for smaller datasets
\end{itemize}

\textbf{K-Means:}

\begin{itemize}
\item
  More scalable and faster
\item
  Requires specifying k
\item
  Better for production deployment
\end{itemize}

\textbf{Both methods converged:} Strong evidence of genuine structure.
\end{frame}

\begin{frame}{Business Value}
\textbf{Enables:}

\begin{itemize}
\item
  Targeted marketing campaigns by segment
\item
  Optimized resource allocation to high-value customers
\item
  Opportunities for customer retention and growth
\item
  Data-driven customer understanding
\end{itemize}

\textbf{ROI:} Improved marketing efficiency and customer lifetime value.
\end{frame}

\begin{frame}{Next Steps}
\textbf{Implementation:}

\begin{enumerate}
\item
  Validate cluster assignments with domain experts
\item
  Implement targeted campaigns for each segment
\item
  Monitor segment-specific KPIs (conversion, AOV, retention)
\item
  Re-run clustering periodically to detect behavioral shifts
\item
  Consider additional features (geographic, demographic) for refinement
\end{enumerate}
\end{frame}

\begin{frame}{Validation in Practice}
\textbf{Important Note:} In real-world unsupervised learning, true
labels do not exist.

\textbf{Cluster validation relies on:}

\begin{itemize}
\item
  Domain expertise
\item
  Business metrics
\item
  Silhouette analysis
\item
  Stability across different methods
\end{itemize}

This synthetic dataset included true labels only for educational
validation purposes.
\end{frame}

\begin{frame}[fragile]{Python Code Resources}
\textbf{Complete implementation available in:}

\begin{itemize}
\item
  \texttt{customer\_clustering\_analysis.ipynb}
\item
  \texttt{fetch\_customer\_data.py} (data generation)
\item
  \texttt{CUSTOMER\_DATA\_DICTIONARY.md} (variable descriptions)
\end{itemize}

\textbf{Key Libraries:}

\begin{itemize}
\item
  scikit-learn (KMeans, StandardScaler)
\item
  scipy (hierarchical clustering)
\item
  pandas, numpy (data manipulation)
\item
  matplotlib, seaborn (visualization)
\end{itemize}
\end{frame}

\begin{frame}{Questions?}
\textbf{ Thank you for your attention! }

Juliho Castillo Colmenares\\
julihocc@tec

Office: Tec de Monterrey CCM Office 1540\\
Office Hours: Monday-Friday, 9:00 AM - 5:00 PM
\end{frame}

\begin{frame}{Next Steps: This Week}
\textbf{For This Week:}

\begin{itemize}
\item
  Review lecture notes thoroughly
\item
  Practice with provided examples
\item
  Complete practice questions
\item
  Prepare for E06 quiz
\end{itemize}
\end{frame}

\begin{frame}{Next Steps: Preparation for Evaluation}
\textbf{Preparation for Evaluation:}

\begin{itemize}
\item
  Understand distance measures and when to use each
\item
  Know linkage methods and their properties
\item
  Practice interpreting dendrograms
\item
  Understand k-means algorithm and convergence
\item
  Be able to explain validation methods
\end{itemize}
\end{frame}

\end{document}




